\section{Gobernanza de seguridad y proceso de revisión}
\subsection{Diseño de la cadena de auditoría}
Cada comando, diff de archivos y resultado de pruebas que genera OpenHands se envía a la cadena de auditoría, y el servicio de auditoría los escribe en el middleware de eventos. Esto permite que el equipo revise fácilmente, después de los hechos, «quién hizo que el agente ejecutara qué comando», y lo contraste con la política de seguridad.

\begin{table}[H]
\centering
\begin{tabular}{p{0.25\textwidth} p{0.4\textwidth} p{0.33\textwidth}}
\toprule
Fase de auditoría & Contenido & Salida \\
\midrule
Captura de eventos & Comandos, navegación, operaciones de archivos & ID de log + marca de tiempo \\
\midrule
Alineación con la política & Verifica si se activan recursos sensibles & Warning / Approval request \\
\midrule
Archivo del informe & Genera correo/issue & Registro de auditoría reproducible \\
\bottomrule
\end{tabular}
\caption{Las tres fases del proceso de auditoría de seguridad.}
\end{table}

\begin{importantbox}{Elementos de la auditoría}
El registro debe incluir el contexto (issue, diff, pruebas), para que una persona pueda juzgar el riesgo con rapidez.
\end{importantbox}

\subsection{Aprobación y permisos}
Cuando el agente planea ejecutar un comando sensible (por ejemplo, una migración de base de datos o un despliegue en la nube), el módulo de aprobación se activa automáticamente. La decisión de aprobación incluye el comando que la desencadena, el entorno de invocación, el resultado esperado, el diff relacionado y el estado de la prueba anterior; quien aprueba solo necesita confirmar o rechazar.

\begin{warningbox}{El ritmo de la aprobación}
No se debe aprobar en exceso al punto de bloquear la iteración; pero en los escenarios de alto riesgo es necesario establecer multi-party approval, para evitar un punto único de falla.
\end{warningbox}

\subsection{Resumen de la sección}
\begin{itemize}
    \item La cadena de auditoría de OpenHands registra el comportamiento del agente, para usarse en revisiones de seguridad posteriores.
    \item Establecer la aprobación y los permisos de manera adecuada puede reforzar la seguridad sin reducir la eficiencia.
\end{itemize}
